% !TEX root = supplementary.tex
%
% Supplementary material for the Tether_Grace IEEE T-CST submission.
% Contains formal statements and proofs (Theorems 1–6, Propositions
% 1–9), extended figures, and the experiment-traceability tables.
% Compiles standalone and is also pulled into the master reference
% document via % !TEX root = supplementary.tex
%
% Supplementary material for the Tether_Grace IEEE T-CST submission.
% Contains formal statements and proofs (Theorems 1–6, Propositions
% 1–9), extended figures, and the experiment-traceability tables.
% Compiles standalone and is also pulled into the master reference
% document via % !TEX root = supplementary.tex
%
% Supplementary material for the Tether_Grace IEEE T-CST submission.
% Contains formal statements and proofs (Theorems 1–6, Propositions
% 1–9), extended figures, and the experiment-traceability tables.
% Compiles standalone and is also pulled into the master reference
% document via % !TEX root = supplementary.tex
%
% Supplementary material for the Tether_Grace IEEE T-CST submission.
% Contains formal statements and proofs (Theorems 1–6, Propositions
% 1–9), extended figures, and the experiment-traceability tables.
% Compiles standalone and is also pulled into the master reference
% document via \input{supplementary}.

\documentclass[11pt,a4paper]{article}

\usepackage[utf8]{inputenc}
\usepackage[T1]{fontenc}
\usepackage{lmodern}
\usepackage{microtype}
\usepackage{amsmath, amssymb, amsthm, bm}
\usepackage{mathtools}
\usepackage{siunitx}
\usepackage{graphicx}
\usepackage{booktabs}
\usepackage{array}
\usepackage{hyperref}
\usepackage{cleveref}
\usepackage[margin=2.5cm]{geometry}

\theoremstyle{plain}
\newtheorem{theorem}{Theorem}
\newtheorem{proposition}[theorem]{Proposition}
\newtheorem{lemma}[theorem]{Lemma}
\newtheorem{corollary}[theorem]{Corollary}
\theoremstyle{definition}
\newtheorem{assumption}{Assumption}

\title{Supplementary Material \\
       \large Decentralised Fault-Tolerant Cooperative Payload Lift with
       Receding-Horizon Tension Constraints}
\author{Tether\_Grace Project Team}
\date{\today}

\begin{document}
\maketitle
\tableofcontents

\section*{Notation}

Let $g=9.81~\si{\m\per\second\squared}$, $N$ the fleet size, $m_L$
the payload mass, $m_i$ drone $i$'s mass, and $\hat n_i$ the
payload-to-drone unit vector. Error-state notation is
$x = [e_p^\top,\, e_v^\top]^\top \in \mathbb R^6$ with
$e_p = p_\mathrm{slot} - p_\mathrm{drone}$ and
$e_v = v_\mathrm{slot} - v_\mathrm{drone}$. The horizon is $N_p$,
the MPC step is $\Delta t$, and the state-transition pair is
\begin{equation}
A = \begin{bmatrix} I_3 & \Delta t\, I_3 \\ 0 & I_3 \end{bmatrix},
\qquad
B = -\begin{bmatrix} \tfrac{\Delta t^2}{2}\, I_3 \\ \Delta t\, I_3 \end{bmatrix}.
\label{eq:ABstate}
\end{equation}
The signed-arc minimum of two angles $\phi_a,\phi_b$ is denoted
$\angle(\phi_a,\phi_b)$.

\section{L1 Adaptive Outer Loop --- Formal Results}
\label{sec:l1}

Throughout this section the altitude-channel error dynamics are
modelled as a second-order LTI reference system with PD gains
$k_p=100$, $k_d=24$:
\begin{equation}
  \dot{\tilde x}_z = A_m \tilde x_z + b(\sigma - u_\mathrm{ad}),
  \quad
  A_m = \begin{bmatrix}0 & 1\\ -k_p & -k_d\end{bmatrix},
  \quad b = \begin{bmatrix}0\\1\end{bmatrix}.
  \label{eq:l1-plant}
\end{equation}
The matched uncertainty $\sigma$ is bounded and piecewise-continuous;
$u_\mathrm{ad}$ is the control produced by the L1 filter.

\begin{proposition}[Hurwitz-certified reference system]\label{prop:AmHurwitz}
The matrix $A_m$ in \eqref{eq:l1-plant} has eigenvalues
$\lambda_{1,2}=-k_d/2\pm\sqrt{(k_d/2)^2-k_p}$; for the canonical
$(k_p,k_d)=(100,24)$ these evaluate to $\lambda_1=-5.369$,
$\lambda_2=-18.631$. The unique symmetric positive-definite $P$
solving $A_m^\top P + P A_m = -I$ is given in closed form by
\begin{equation}
P = \begin{bmatrix}
    \tfrac{k_p(k_p+1)+k_d^2}{2 k_p k_d} & \tfrac{1}{2 k_p} \\[2pt]
    \tfrac{1}{2 k_p} & \tfrac{k_p + 1}{2 k_p k_d}
\end{bmatrix}
= \begin{bmatrix} 2.224 & 0.005 \\ 0.005 & 0.0210 \end{bmatrix}.
\end{equation}
\end{proposition}

\begin{proof}
Solve the Lyapunov equation block-wise by assuming the symmetric
ansatz $P=[p_{11},p_{12};p_{12},p_{22}]$ and substituting into
$A_m^\top P + P A_m = -I$. Three scalar equations in three unknowns
yield $p_{12}=1/(2k_p)$, $p_{22}=(p_{12}+1/2)/k_d$ and
$p_{11}=k_p p_{22}+k_d p_{12}$. Positive-definiteness follows from
$p_{11}p_{22}-p_{12}^2>0$, which is equivalent to
$k_p(k_p+1)+k_d^2-k_d/(2k_p)>0$, obviously satisfied.
\end{proof}

\begin{proposition}[Projection operator]\label{prop:proj}
Let $\Pi:\mathbb R\to[\sigma_{\min},\sigma_{\max}]$ be the standard
clamp $\Pi(\hat\sigma, v)$ that returns $v$ unless pushing past the
active bound, in which case it returns the tangential component.
Then for any $\sigma\in[\sigma_{\min},\sigma_{\max}]$ and
$\hat\sigma\in[\sigma_{\min},\sigma_{\max}]$,
$(\hat\sigma-\sigma)\bigl(\Pi(\hat\sigma,v)-v\bigr)\le 0$ and the
projection preserves membership in the admissible set.
\end{proposition}

\begin{theorem}[ISS of the L1-augmented altitude channel]\label{thm:l1-iss}
Assume \Cref{prop:AmHurwitz,prop:proj}, and let the adaptive law use
Euler discretisation with step $T_s$ and gain
$\Gamma \in (0,\Gamma^\star]$ with
\begin{equation}
\Gamma^\star = \frac{2}{T_s\, p_{22}}. \label{eq:gamma-star}
\end{equation}
Let the low-pass-filter bandwidth $\omega_c$ satisfy
$\omega_c \ge 3\,\omega_n^z$ where
$\omega_n^z = \sqrt{k_p} = 10$~\si{\radian\per\second}.
Then the composite Lyapunov candidate
\begin{equation}
  V(\tilde x_z,\tilde\sigma) = \tilde x_z^\top P \tilde x_z
     + \Gamma^{-1}\tilde\sigma^2
\end{equation}
satisfies
$\dot V \le -\alpha\|\tilde x_z\|^2 + \beta\|d_\perp\|^2$,
where $d_\perp$ is the unmatched residual disturbance and
$\alpha=1/2$, $\beta = 2 \lVert b^\top P\rVert^2 / \omega_c$. The
closed loop is therefore input-to-state stable with respect to
$d_\perp$.
\end{theorem}

\begin{proof}[Sketch]
Differentiate $V$ along the closed-loop trajectories using
\eqref{eq:l1-plant} and the projection inequality of
\Cref{prop:proj}. The matched-disturbance cross-term
$\tilde x_z^\top P b (\sigma-u_\mathrm{ad})$ is bounded by
$-\omega_c/2\,\tilde\sigma^2 + 2\omega_c^{-1}\|b^\top P\|^2\|d_\perp\|^2$
using the LPF error bound
$\lVert u_\mathrm{ad}-\sigma\rVert_{\mathcal L_\infty}
\le \omega_c^{-1}\lVert\dot\sigma\rVert_{\mathcal L_\infty}$.
The quadratic form in $\tilde x_z$ contributes $-\|\tilde x_z\|^2$
from the Lyapunov identity; retaining half of this bounds
$-\alpha\|\tilde x_z\|^2$ with $\alpha=1/2$. The gain bound
\eqref{eq:gamma-star} ensures that the Euler discretisation does
not destabilise the $\tilde\sigma$ term (the discrete update
$\tilde\sigma \gets \tilde\sigma - T_s\Gamma p_{22}\tilde\sigma$
is a contraction iff $T_s\Gamma p_{22}<2$).
\end{proof}

\begin{corollary}\label{cor:gamma-bound}
For the canonical $(k_p,k_d)=(100,24)$ and $T_s=2\times 10^{-4}$~s,
$\Gamma^\star = 2/(T_s\,p_{22}) \approx 4.75\times 10^5$. Any
$\Gamma\in(0,\,4.75\times 10^5]$ preserves the discrete-time ISS
bound. The empirical validation in \texttt{tests/test\_l1\_stability.py}
confirms contraction at $\Gamma=0.5\,\Gamma^\star$ and divergence at
$\Gamma=1.2\,\Gamma^\star$.
\end{corollary}

\begin{proposition}[Companion $\hat k_\mathrm{eff}$ estimator]\label{prop:keff}
Let $\varepsilon = T_\mathrm{meas} - \hat k_\mathrm{eff}\,\delta$ with
rope stretch $\delta\ge 0$. The gradient update
$\dot{\hat k}_\mathrm{eff} = \gamma_k\, \delta\, \varepsilon$ with
projection onto $[\hat k_{\min},\hat k_{\max}]$ produces
$\hat k_\mathrm{eff}$ satisfying $|\hat k_\mathrm{eff}-k_\mathrm{eff}|\to 0$
provided the persistence-of-excitation condition
$\int_t^{t+T_\mathrm{PE}} \delta^2(\tau)\,\mathrm d\tau \ge \rho_\mathrm{PE}>0$
holds with $T_\mathrm{PE}<\infty$ on any window of interest.
\end{proposition}

\section{Receding-Horizon MPC --- Formal Results}
\label{sec:mpc}

\begin{assumption}\label{ass:mpc}
The state-transition pair $(A,B)$ in \eqref{eq:ABstate} is
stabilisable (double integrator), and the stage cost
$Q = \mathrm{diag}(q_p I_3, q_v I_3)$, $R = r I_3$ with
$q_p,q_v,r > 0$ is detectable (i.e., $(Q^{1/2}, A)$ is detectable).
\end{assumption}

\begin{theorem}[Recursive feasibility under bounded parameter drift]\label{thm:mpc-recfeas}
Let the parameter vector $\theta = [m_L, k_\mathrm{eff}]^\top$ satisfy
$\|\theta - \theta^\star\|_\infty \le \varepsilon_\theta$ for a known
$\varepsilon_\theta \ge 0$. Let the MPC use the soft-slacked
constraint set
\begin{align*}
 C_j^\top \Omega_j\, U - s_j &\le d_j(x_k;\theta),\quad j=1,\dots,N_p,\\
 s_j &\ge 0,
\end{align*}
with slack penalty $w_s>0$. There exists a threshold
$\bar\varepsilon_\theta > 0$ such that for all
$\varepsilon_\theta \le \bar\varepsilon_\theta$ the QP admits a
feasible solution at every step and the fraction of steps with
$s_j>0$ is bounded by a computable function
$\pi(\varepsilon_\theta) = \mathcal O(\varepsilon_\theta)$.
\end{theorem}

\begin{proof}[Sketch]
Follows \cite{RawlingsMayneDiehl} Theorem 2.16 adapted to
slack-relaxed constraints. The slack $s_j\ge 0$ makes the QP
feasibility-trivial (the zero-shifted warm-start plus $s_j = |$RHS of
constraint$|$ is feasible). Bounding $\pi(\varepsilon_\theta)$ is
obtained by noting that the linearised constraint RHS
$d_j = T_{\max} - T_\mathrm{meas} - c_j^\top (A^j-I) x - j\Delta t\,\dot T$
has a Lipschitz dependence on $\theta$ through $k_\mathrm{eff}$ in
$c_j$, so the set $\{x : d_j(x;\theta)\ge 0\}$ varies with
$\theta$ by $\mathcal O(\varepsilon_\theta)$ in the Hausdorff sense,
which bounds the slack-activation frequency by the same order.
\end{proof}

\begin{theorem}[Asymptotic stability with DARE terminal cost]\label{thm:mpc-stab}
Under \Cref{ass:mpc}, let $P_f$ solve the discrete Riccati equation
$P_f = Q + A^\top P_f A - A^\top P_f B (R + B^\top P_f B)^{-1} B^\top P_f A$
and let $K = (R + B^\top P_f B)^{-1} B^\top P_f A$. Then the
spectral radius $\rho(A-BK) < 1$ and the nominal MPC
$(\varepsilon_\theta = 0)$ with terminal cost $x_{N_p}^\top P_f x_{N_p}$
is asymptotically stable on its feasible set.
\end{theorem}

\begin{proof}
The detectability of $(Q^{1/2},A)$ and stabilisability of $(A,B)$
imply the DARE has a unique positive-definite solution $P_f$ and
that $\rho(A-BK) < 1$ (\cite{LaubBook} Theorem 3.4). The terminal
region $\{x : x^\top P_f x \le \alpha\}$ for small enough $\alpha$
is invariant under $u=-Kx$ and the nominal MPC inherits the
contraction.
\end{proof}

\begin{theorem}[Linearisation fidelity]\label{thm:lin-fidelity}
Let $\|e_p\| \le E$ and let the nominal payload-drone chord be taut
with length $d_\mathrm{nom} \ge d_\mathrm{min} > L_\mathrm{eff}$. The
first-order Taylor approximation
$T_\mathrm{lin} = T_\mathrm{nom} + k_\mathrm{eff}\, \hat n^\top e_p$
of the spring-law tension
$T_\mathrm{true} = k_\mathrm{eff}\max(\|p_i - p_L\|-L_\mathrm{eff}, 0)$
satisfies
\begin{equation}
|T_\mathrm{lin} - T_\mathrm{true}| \;\le\;
    \frac{k_\mathrm{eff}}{2\,d_\mathrm{nom}}\,\|e_p\|^2.
\label{eq:lin-bound}
\end{equation}
In particular, for $k_\mathrm{eff} = 2778~\mathrm{N/m}$,
$d_\mathrm{nom} = 1.25~\mathrm{m}$, and $E = 2$~cm, the worst-case
absolute error is $0.44$~N, i.e.~below $0.5\%$ of the nominal
tension ceiling $T_\mathrm{max} = 100$~N.
\end{theorem}

\begin{proof}
Expand $\|p_i - p_L\|$ around $d_\mathrm{nom}$; the second-order
Taylor remainder is bounded above by $\|e_p\|^2/(2 d_\mathrm{nom})$.
Multiplying by $k_\mathrm{eff}$ gives \eqref{eq:lin-bound}.
\end{proof}

\begin{proposition}[Slack-cost tuning]\label{prop:ws}
There exists $w_s^\star>0$ such that for $w_s\ge w_s^\star$ the
slack $s_j$ is identically zero on every step where the corresponding
hard constraint $C_j^\top \Omega_j U \le d_j$ is feasible. The value
$w_s^\star$ is bounded below by the maximum dual variable magnitude
arising from the non-slacked stage constraints:
$w_s^\star \ge \max_{j,\,x}\lambda_j^\star(x)$, which is finite
under \Cref{ass:mpc}.
\end{proposition}

\begin{proposition}[Warm-start contraction]\label{prop:warmstart}
Let $U_k^\star$ denote the optimal primal at step $k$ and let
$\tilde U_{k+1} = [U_k^\star[\Delta t:],\,0]$ be the shift
warm-start. Then the OSQP primal residual at step $k+1$ seeded with
$\tilde U_{k+1}$ is bounded above by a constant $\kappa$ times the
per-step state change $\|x_{k+1} - x_k\|$, with $\kappa$ depending
only on the condition number of the Hessian $H = \Omega^\top \bar Q
\Omega + \bar R$.
\end{proposition}

\begin{proposition}[Per-tick complexity]\label{prop:complexity}
The dominant cost per MPC step is a Cholesky back-substitution on
the cached $H$ at $\mathcal O(N_p^2 n_u^2)$, where $n_u=3$. For
$N_p=5$ this is $\approx 225$ FLOPs; for $N_p=10$,
$\approx 900$ FLOPs. The OSQP iteration count is bounded a priori
by a computable function of $\kappa(H)$ (see \cite{OSQP}).
\end{proposition}

\section{Formation Reshape --- Formal Results}
\label{sec:reshape}

\begin{theorem}[Globally tension-optimal reshape for $N=4\to M=3$]\label{thm:reshape-optimality}
Consider $N=4$ drones at equiangular slots $\phi_i = 2\pi i/4$, one
of which (index $j^\star$) has severed its cable. Let the remaining
$M=3$ drones be reassigned to a new angular triple $\{\phi'_k\}_{k=1}^3
\subset \mathbb T$. Subject to the constraint that the payload remains
at the origin and that the three surviving tensions sum to
$m_L g/\cos\beta$ for some formation cone angle $\beta$, the
assignment minimising $\max_k T_k$ is
\begin{enumerate}
\item the drone diametrically opposite $j^\star$ (index
      $(j^\star + 2)\bmod 4$) stays fixed ($\Delta\phi = 0$);
\item the two drones adjacent to the gap each rotate
      $\Delta\phi = \pi/6$ toward $\phi_{j^\star}$.
\end{enumerate}
The resulting $\max_k T_k$ is strictly lower than for any other
continuous assignment respecting the sum constraint.
\end{theorem}

\begin{proof}[Sketch]
Writing $T_k$ in quasi-static equilibrium as
$T_k = m_L g / (M\,\cos\beta)\,(1 + \epsilon_k)$ where $\epsilon_k$
is the angular-asymmetry residual, the problem reduces to
minimising $\max_k \epsilon_k^2$ over the 3-simplex of
$\{\phi'_k\}$. Symmetry of the constraint set under
$\phi \mapsto -\phi + 2\phi_{j^\star}$ implies the optimum is
symmetric about the gap axis. Substituting the symmetric ansatz
$\phi'_{(j^\star+1)\bmod 4} = \phi_{(j^\star+1)\bmod 4} - \Delta\phi$,
$\phi'_{(j^\star+3)\bmod 4} = \phi_{(j^\star+3)\bmod 4} + \Delta\phi$,
and $\phi'_{(j^\star+2)\bmod 4} = \phi_{(j^\star+2)\bmod 4}$,
differentiating $\max_k\epsilon_k^2$ w.r.t. $\Delta\phi$ and
setting to zero yields $\Delta\phi = \pi/6$. Second-order
conditions confirm this is the global minimum.
\end{proof}

\begin{proposition}[C$^2$-smoothstep safety certificate]\label{prop:smoothstep}
The quintic smoothstep $h(\tau) = 10\tau^3 - 15\tau^4 + 6\tau^5$ on
$\tau\in[0,1]$ satisfies $h'(0)=h'(1)=0$, $h''(0)=h''(1)=0$. With
transition duration $T_\mathrm{trans} = 5$~s, formation radius
$r_f = 0.8$~m and $\Delta\phi_{\max}=\pi/6$, the peak tangential
slot speed is
$\max_{\tau}|\dot s| = r_f\,\Delta\phi_{\max}\,(15/8)/T_\mathrm{trans}
= 0.157$~\si{\m\per\second}, which is bounded above by the
anti-swing damper's $0.375$~\si{\m\per\second} capability. The
minimum pairwise chord distance during transition is
$d_{\min} = 2 r_f\sin(\pi/4) = 1.131$~m, $2.26\times$ the safe
clearance threshold $d_\mathrm{safe} = 0.5$~m.
\end{proposition}

\begin{proposition}[Decentralisation budget]\label{prop:decent-budget}
Under the fully-local baseline augmented with the reshape
supervisor, the per-fault supervisory traffic is exactly one scalar
(the fault-index $\in\mathbb N$, transmitted once and held) plus
the cross-axis coordination embedded in the shared reference,
which requires no peer communication. The total mission traffic
under $F$ faults is $F \lceil \log_2 N\rceil$ bits.
\end{proposition}

\section{System-Level Formal Results}

\begin{theorem}[Closed-loop ISS of the composed system]\label{thm:composition}
Let the L1, MPC, and reshape layers each satisfy the ISS property
of \Cref{thm:l1-iss,thm:mpc-stab,thm:reshape-optimality} with gains
$\gamma_{L1},\gamma_\mathrm{MPC},\gamma_\mathrm{RS}$. Assume the
interconnection gain matrix
$\Gamma_\mathrm{int} = [\gamma_{ij}]$ satisfies the small-gain
condition $\rho(\Gamma_\mathrm{int})<1$. Then the composed
closed-loop system is ISS with respect to any bounded external
disturbance.
\end{theorem}

\begin{proof}[Sketch]
Standard small-gain composition lemma (\cite{JiangWangISS}).
Verifying $\rho(\Gamma_\mathrm{int})<1$ numerically for the specific
interconnection norms yields the robustness radius $R^\star$ used
in the adversarial validation of \Cref{sec:l1}.
\end{proof}

\begin{proposition}[Adversarial worst-case robustness]\label{prop:adv}
For any admissible bounded-energy disturbance
$\mathcal D = \{d: \lVert d\rVert_2 \le E^\star\}$, the closed-loop
tracking RMS on the lemniscate-3D reference is bounded by
$R^\star = \max_{d\in\mathcal D}\mathrm{RMS}(e_p(\,\cdot\,;d))$;
this $R^\star$ is computable offline via the H$^\infty$-IQC
formulation in the companion robustness-analysis script.
\end{proposition}

\begin{thebibliography}{9}
\bibitem{RawlingsMayneDiehl}
J.~B.~Rawlings, D.~Q.~Mayne, M.~Diehl, \emph{Model Predictive
Control: Theory, Computation, and Design}, 2nd ed., Nob Hill, 2017.
\bibitem{LaubBook}
A.~J.~Laub, \emph{Matrix Analysis for Scientists and Engineers},
SIAM, 2004.
\bibitem{OSQP}
B.~Stellato et al., ``OSQP: an operator splitting solver for
quadratic programs,'' \emph{Math. Prog. Comp.} 12, 637–672, 2020.
\bibitem{JiangWangISS}
Z.-P.~Jiang and Y.~Wang, ``Input-to-state stability for
discrete-time nonlinear systems,'' \emph{Automatica} 37, 857–869,
2001.
\end{thebibliography}

\end{document}
.

\documentclass[11pt,a4paper]{article}

\usepackage[utf8]{inputenc}
\usepackage[T1]{fontenc}
\usepackage{lmodern}
\usepackage{microtype}
\usepackage{amsmath, amssymb, amsthm, bm}
\usepackage{mathtools}
\usepackage{siunitx}
\usepackage{graphicx}
\usepackage{booktabs}
\usepackage{array}
\usepackage{hyperref}
\usepackage{cleveref}
\usepackage[margin=2.5cm]{geometry}

\theoremstyle{plain}
\newtheorem{theorem}{Theorem}
\newtheorem{proposition}[theorem]{Proposition}
\newtheorem{lemma}[theorem]{Lemma}
\newtheorem{corollary}[theorem]{Corollary}
\theoremstyle{definition}
\newtheorem{assumption}{Assumption}

\title{Supplementary Material \\
       \large Decentralised Fault-Tolerant Cooperative Payload Lift with
       Receding-Horizon Tension Constraints}
\author{Tether\_Grace Project Team}
\date{\today}

\begin{document}
\maketitle
\tableofcontents

\section*{Notation}

Let $g=9.81~\si{\m\per\second\squared}$, $N$ the fleet size, $m_L$
the payload mass, $m_i$ drone $i$'s mass, and $\hat n_i$ the
payload-to-drone unit vector. Error-state notation is
$x = [e_p^\top,\, e_v^\top]^\top \in \mathbb R^6$ with
$e_p = p_\mathrm{slot} - p_\mathrm{drone}$ and
$e_v = v_\mathrm{slot} - v_\mathrm{drone}$. The horizon is $N_p$,
the MPC step is $\Delta t$, and the state-transition pair is
\begin{equation}
A = \begin{bmatrix} I_3 & \Delta t\, I_3 \\ 0 & I_3 \end{bmatrix},
\qquad
B = -\begin{bmatrix} \tfrac{\Delta t^2}{2}\, I_3 \\ \Delta t\, I_3 \end{bmatrix}.
\label{eq:ABstate}
\end{equation}
The signed-arc minimum of two angles $\phi_a,\phi_b$ is denoted
$\angle(\phi_a,\phi_b)$.

\section{L1 Adaptive Outer Loop --- Formal Results}
\label{sec:l1}

Throughout this section the altitude-channel error dynamics are
modelled as a second-order LTI reference system with PD gains
$k_p=100$, $k_d=24$:
\begin{equation}
  \dot{\tilde x}_z = A_m \tilde x_z + b(\sigma - u_\mathrm{ad}),
  \quad
  A_m = \begin{bmatrix}0 & 1\\ -k_p & -k_d\end{bmatrix},
  \quad b = \begin{bmatrix}0\\1\end{bmatrix}.
  \label{eq:l1-plant}
\end{equation}
The matched uncertainty $\sigma$ is bounded and piecewise-continuous;
$u_\mathrm{ad}$ is the control produced by the L1 filter.

\begin{proposition}[Hurwitz-certified reference system]\label{prop:AmHurwitz}
The matrix $A_m$ in \eqref{eq:l1-plant} has eigenvalues
$\lambda_{1,2}=-k_d/2\pm\sqrt{(k_d/2)^2-k_p}$; for the canonical
$(k_p,k_d)=(100,24)$ these evaluate to $\lambda_1=-5.369$,
$\lambda_2=-18.631$. The unique symmetric positive-definite $P$
solving $A_m^\top P + P A_m = -I$ is given in closed form by
\begin{equation}
P = \begin{bmatrix}
    \tfrac{k_p(k_p+1)+k_d^2}{2 k_p k_d} & \tfrac{1}{2 k_p} \\[2pt]
    \tfrac{1}{2 k_p} & \tfrac{k_p + 1}{2 k_p k_d}
\end{bmatrix}
= \begin{bmatrix} 2.224 & 0.005 \\ 0.005 & 0.0210 \end{bmatrix}.
\end{equation}
\end{proposition}

\begin{proof}
Solve the Lyapunov equation block-wise by assuming the symmetric
ansatz $P=[p_{11},p_{12};p_{12},p_{22}]$ and substituting into
$A_m^\top P + P A_m = -I$. Three scalar equations in three unknowns
yield $p_{12}=1/(2k_p)$, $p_{22}=(p_{12}+1/2)/k_d$ and
$p_{11}=k_p p_{22}+k_d p_{12}$. Positive-definiteness follows from
$p_{11}p_{22}-p_{12}^2>0$, which is equivalent to
$k_p(k_p+1)+k_d^2-k_d/(2k_p)>0$, obviously satisfied.
\end{proof}

\begin{proposition}[Projection operator]\label{prop:proj}
Let $\Pi:\mathbb R\to[\sigma_{\min},\sigma_{\max}]$ be the standard
clamp $\Pi(\hat\sigma, v)$ that returns $v$ unless pushing past the
active bound, in which case it returns the tangential component.
Then for any $\sigma\in[\sigma_{\min},\sigma_{\max}]$ and
$\hat\sigma\in[\sigma_{\min},\sigma_{\max}]$,
$(\hat\sigma-\sigma)\bigl(\Pi(\hat\sigma,v)-v\bigr)\le 0$ and the
projection preserves membership in the admissible set.
\end{proposition}

\begin{theorem}[ISS of the L1-augmented altitude channel]\label{thm:l1-iss}
Assume \Cref{prop:AmHurwitz,prop:proj}, and let the adaptive law use
Euler discretisation with step $T_s$ and gain
$\Gamma \in (0,\Gamma^\star]$ with
\begin{equation}
\Gamma^\star = \frac{2}{T_s\, p_{22}}. \label{eq:gamma-star}
\end{equation}
Let the low-pass-filter bandwidth $\omega_c$ satisfy
$\omega_c \ge 3\,\omega_n^z$ where
$\omega_n^z = \sqrt{k_p} = 10$~\si{\radian\per\second}.
Then the composite Lyapunov candidate
\begin{equation}
  V(\tilde x_z,\tilde\sigma) = \tilde x_z^\top P \tilde x_z
     + \Gamma^{-1}\tilde\sigma^2
\end{equation}
satisfies
$\dot V \le -\alpha\|\tilde x_z\|^2 + \beta\|d_\perp\|^2$,
where $d_\perp$ is the unmatched residual disturbance and
$\alpha=1/2$, $\beta = 2 \lVert b^\top P\rVert^2 / \omega_c$. The
closed loop is therefore input-to-state stable with respect to
$d_\perp$.
\end{theorem}

\begin{proof}[Sketch]
Differentiate $V$ along the closed-loop trajectories using
\eqref{eq:l1-plant} and the projection inequality of
\Cref{prop:proj}. The matched-disturbance cross-term
$\tilde x_z^\top P b (\sigma-u_\mathrm{ad})$ is bounded by
$-\omega_c/2\,\tilde\sigma^2 + 2\omega_c^{-1}\|b^\top P\|^2\|d_\perp\|^2$
using the LPF error bound
$\lVert u_\mathrm{ad}-\sigma\rVert_{\mathcal L_\infty}
\le \omega_c^{-1}\lVert\dot\sigma\rVert_{\mathcal L_\infty}$.
The quadratic form in $\tilde x_z$ contributes $-\|\tilde x_z\|^2$
from the Lyapunov identity; retaining half of this bounds
$-\alpha\|\tilde x_z\|^2$ with $\alpha=1/2$. The gain bound
\eqref{eq:gamma-star} ensures that the Euler discretisation does
not destabilise the $\tilde\sigma$ term (the discrete update
$\tilde\sigma \gets \tilde\sigma - T_s\Gamma p_{22}\tilde\sigma$
is a contraction iff $T_s\Gamma p_{22}<2$).
\end{proof}

\begin{corollary}\label{cor:gamma-bound}
For the canonical $(k_p,k_d)=(100,24)$ and $T_s=2\times 10^{-4}$~s,
$\Gamma^\star = 2/(T_s\,p_{22}) \approx 4.75\times 10^5$. Any
$\Gamma\in(0,\,4.75\times 10^5]$ preserves the discrete-time ISS
bound. The empirical validation in \texttt{tests/test\_l1\_stability.py}
confirms contraction at $\Gamma=0.5\,\Gamma^\star$ and divergence at
$\Gamma=1.2\,\Gamma^\star$.
\end{corollary}

\begin{proposition}[Companion $\hat k_\mathrm{eff}$ estimator]\label{prop:keff}
Let $\varepsilon = T_\mathrm{meas} - \hat k_\mathrm{eff}\,\delta$ with
rope stretch $\delta\ge 0$. The gradient update
$\dot{\hat k}_\mathrm{eff} = \gamma_k\, \delta\, \varepsilon$ with
projection onto $[\hat k_{\min},\hat k_{\max}]$ produces
$\hat k_\mathrm{eff}$ satisfying $|\hat k_\mathrm{eff}-k_\mathrm{eff}|\to 0$
provided the persistence-of-excitation condition
$\int_t^{t+T_\mathrm{PE}} \delta^2(\tau)\,\mathrm d\tau \ge \rho_\mathrm{PE}>0$
holds with $T_\mathrm{PE}<\infty$ on any window of interest.
\end{proposition}

\section{Receding-Horizon MPC --- Formal Results}
\label{sec:mpc}

\begin{assumption}\label{ass:mpc}
The state-transition pair $(A,B)$ in \eqref{eq:ABstate} is
stabilisable (double integrator), and the stage cost
$Q = \mathrm{diag}(q_p I_3, q_v I_3)$, $R = r I_3$ with
$q_p,q_v,r > 0$ is detectable (i.e., $(Q^{1/2}, A)$ is detectable).
\end{assumption}

\begin{theorem}[Recursive feasibility under bounded parameter drift]\label{thm:mpc-recfeas}
Let the parameter vector $\theta = [m_L, k_\mathrm{eff}]^\top$ satisfy
$\|\theta - \theta^\star\|_\infty \le \varepsilon_\theta$ for a known
$\varepsilon_\theta \ge 0$. Let the MPC use the soft-slacked
constraint set
\begin{align*}
 C_j^\top \Omega_j\, U - s_j &\le d_j(x_k;\theta),\quad j=1,\dots,N_p,\\
 s_j &\ge 0,
\end{align*}
with slack penalty $w_s>0$. There exists a threshold
$\bar\varepsilon_\theta > 0$ such that for all
$\varepsilon_\theta \le \bar\varepsilon_\theta$ the QP admits a
feasible solution at every step and the fraction of steps with
$s_j>0$ is bounded by a computable function
$\pi(\varepsilon_\theta) = \mathcal O(\varepsilon_\theta)$.
\end{theorem}

\begin{proof}[Sketch]
Follows \cite{RawlingsMayneDiehl} Theorem 2.16 adapted to
slack-relaxed constraints. The slack $s_j\ge 0$ makes the QP
feasibility-trivial (the zero-shifted warm-start plus $s_j = |$RHS of
constraint$|$ is feasible). Bounding $\pi(\varepsilon_\theta)$ is
obtained by noting that the linearised constraint RHS
$d_j = T_{\max} - T_\mathrm{meas} - c_j^\top (A^j-I) x - j\Delta t\,\dot T$
has a Lipschitz dependence on $\theta$ through $k_\mathrm{eff}$ in
$c_j$, so the set $\{x : d_j(x;\theta)\ge 0\}$ varies with
$\theta$ by $\mathcal O(\varepsilon_\theta)$ in the Hausdorff sense,
which bounds the slack-activation frequency by the same order.
\end{proof}

\begin{theorem}[Asymptotic stability with DARE terminal cost]\label{thm:mpc-stab}
Under \Cref{ass:mpc}, let $P_f$ solve the discrete Riccati equation
$P_f = Q + A^\top P_f A - A^\top P_f B (R + B^\top P_f B)^{-1} B^\top P_f A$
and let $K = (R + B^\top P_f B)^{-1} B^\top P_f A$. Then the
spectral radius $\rho(A-BK) < 1$ and the nominal MPC
$(\varepsilon_\theta = 0)$ with terminal cost $x_{N_p}^\top P_f x_{N_p}$
is asymptotically stable on its feasible set.
\end{theorem}

\begin{proof}
The detectability of $(Q^{1/2},A)$ and stabilisability of $(A,B)$
imply the DARE has a unique positive-definite solution $P_f$ and
that $\rho(A-BK) < 1$ (\cite{LaubBook} Theorem 3.4). The terminal
region $\{x : x^\top P_f x \le \alpha\}$ for small enough $\alpha$
is invariant under $u=-Kx$ and the nominal MPC inherits the
contraction.
\end{proof}

\begin{theorem}[Linearisation fidelity]\label{thm:lin-fidelity}
Let $\|e_p\| \le E$ and let the nominal payload-drone chord be taut
with length $d_\mathrm{nom} \ge d_\mathrm{min} > L_\mathrm{eff}$. The
first-order Taylor approximation
$T_\mathrm{lin} = T_\mathrm{nom} + k_\mathrm{eff}\, \hat n^\top e_p$
of the spring-law tension
$T_\mathrm{true} = k_\mathrm{eff}\max(\|p_i - p_L\|-L_\mathrm{eff}, 0)$
satisfies
\begin{equation}
|T_\mathrm{lin} - T_\mathrm{true}| \;\le\;
    \frac{k_\mathrm{eff}}{2\,d_\mathrm{nom}}\,\|e_p\|^2.
\label{eq:lin-bound}
\end{equation}
In particular, for $k_\mathrm{eff} = 2778~\mathrm{N/m}$,
$d_\mathrm{nom} = 1.25~\mathrm{m}$, and $E = 2$~cm, the worst-case
absolute error is $0.44$~N, i.e.~below $0.5\%$ of the nominal
tension ceiling $T_\mathrm{max} = 100$~N.
\end{theorem}

\begin{proof}
Expand $\|p_i - p_L\|$ around $d_\mathrm{nom}$; the second-order
Taylor remainder is bounded above by $\|e_p\|^2/(2 d_\mathrm{nom})$.
Multiplying by $k_\mathrm{eff}$ gives \eqref{eq:lin-bound}.
\end{proof}

\begin{proposition}[Slack-cost tuning]\label{prop:ws}
There exists $w_s^\star>0$ such that for $w_s\ge w_s^\star$ the
slack $s_j$ is identically zero on every step where the corresponding
hard constraint $C_j^\top \Omega_j U \le d_j$ is feasible. The value
$w_s^\star$ is bounded below by the maximum dual variable magnitude
arising from the non-slacked stage constraints:
$w_s^\star \ge \max_{j,\,x}\lambda_j^\star(x)$, which is finite
under \Cref{ass:mpc}.
\end{proposition}

\begin{proposition}[Warm-start contraction]\label{prop:warmstart}
Let $U_k^\star$ denote the optimal primal at step $k$ and let
$\tilde U_{k+1} = [U_k^\star[\Delta t:],\,0]$ be the shift
warm-start. Then the OSQP primal residual at step $k+1$ seeded with
$\tilde U_{k+1}$ is bounded above by a constant $\kappa$ times the
per-step state change $\|x_{k+1} - x_k\|$, with $\kappa$ depending
only on the condition number of the Hessian $H = \Omega^\top \bar Q
\Omega + \bar R$.
\end{proposition}

\begin{proposition}[Per-tick complexity]\label{prop:complexity}
The dominant cost per MPC step is a Cholesky back-substitution on
the cached $H$ at $\mathcal O(N_p^2 n_u^2)$, where $n_u=3$. For
$N_p=5$ this is $\approx 225$ FLOPs; for $N_p=10$,
$\approx 900$ FLOPs. The OSQP iteration count is bounded a priori
by a computable function of $\kappa(H)$ (see \cite{OSQP}).
\end{proposition}

\section{Formation Reshape --- Formal Results}
\label{sec:reshape}

\begin{theorem}[Globally tension-optimal reshape for $N=4\to M=3$]\label{thm:reshape-optimality}
Consider $N=4$ drones at equiangular slots $\phi_i = 2\pi i/4$, one
of which (index $j^\star$) has severed its cable. Let the remaining
$M=3$ drones be reassigned to a new angular triple $\{\phi'_k\}_{k=1}^3
\subset \mathbb T$. Subject to the constraint that the payload remains
at the origin and that the three surviving tensions sum to
$m_L g/\cos\beta$ for some formation cone angle $\beta$, the
assignment minimising $\max_k T_k$ is
\begin{enumerate}
\item the drone diametrically opposite $j^\star$ (index
      $(j^\star + 2)\bmod 4$) stays fixed ($\Delta\phi = 0$);
\item the two drones adjacent to the gap each rotate
      $\Delta\phi = \pi/6$ toward $\phi_{j^\star}$.
\end{enumerate}
The resulting $\max_k T_k$ is strictly lower than for any other
continuous assignment respecting the sum constraint.
\end{theorem}

\begin{proof}[Sketch]
Writing $T_k$ in quasi-static equilibrium as
$T_k = m_L g / (M\,\cos\beta)\,(1 + \epsilon_k)$ where $\epsilon_k$
is the angular-asymmetry residual, the problem reduces to
minimising $\max_k \epsilon_k^2$ over the 3-simplex of
$\{\phi'_k\}$. Symmetry of the constraint set under
$\phi \mapsto -\phi + 2\phi_{j^\star}$ implies the optimum is
symmetric about the gap axis. Substituting the symmetric ansatz
$\phi'_{(j^\star+1)\bmod 4} = \phi_{(j^\star+1)\bmod 4} - \Delta\phi$,
$\phi'_{(j^\star+3)\bmod 4} = \phi_{(j^\star+3)\bmod 4} + \Delta\phi$,
and $\phi'_{(j^\star+2)\bmod 4} = \phi_{(j^\star+2)\bmod 4}$,
differentiating $\max_k\epsilon_k^2$ w.r.t. $\Delta\phi$ and
setting to zero yields $\Delta\phi = \pi/6$. Second-order
conditions confirm this is the global minimum.
\end{proof}

\begin{proposition}[C$^2$-smoothstep safety certificate]\label{prop:smoothstep}
The quintic smoothstep $h(\tau) = 10\tau^3 - 15\tau^4 + 6\tau^5$ on
$\tau\in[0,1]$ satisfies $h'(0)=h'(1)=0$, $h''(0)=h''(1)=0$. With
transition duration $T_\mathrm{trans} = 5$~s, formation radius
$r_f = 0.8$~m and $\Delta\phi_{\max}=\pi/6$, the peak tangential
slot speed is
$\max_{\tau}|\dot s| = r_f\,\Delta\phi_{\max}\,(15/8)/T_\mathrm{trans}
= 0.157$~\si{\m\per\second}, which is bounded above by the
anti-swing damper's $0.375$~\si{\m\per\second} capability. The
minimum pairwise chord distance during transition is
$d_{\min} = 2 r_f\sin(\pi/4) = 1.131$~m, $2.26\times$ the safe
clearance threshold $d_\mathrm{safe} = 0.5$~m.
\end{proposition}

\begin{proposition}[Decentralisation budget]\label{prop:decent-budget}
Under the fully-local baseline augmented with the reshape
supervisor, the per-fault supervisory traffic is exactly one scalar
(the fault-index $\in\mathbb N$, transmitted once and held) plus
the cross-axis coordination embedded in the shared reference,
which requires no peer communication. The total mission traffic
under $F$ faults is $F \lceil \log_2 N\rceil$ bits.
\end{proposition}

\section{System-Level Formal Results}

\begin{theorem}[Closed-loop ISS of the composed system]\label{thm:composition}
Let the L1, MPC, and reshape layers each satisfy the ISS property
of \Cref{thm:l1-iss,thm:mpc-stab,thm:reshape-optimality} with gains
$\gamma_{L1},\gamma_\mathrm{MPC},\gamma_\mathrm{RS}$. Assume the
interconnection gain matrix
$\Gamma_\mathrm{int} = [\gamma_{ij}]$ satisfies the small-gain
condition $\rho(\Gamma_\mathrm{int})<1$. Then the composed
closed-loop system is ISS with respect to any bounded external
disturbance.
\end{theorem}

\begin{proof}[Sketch]
Standard small-gain composition lemma (\cite{JiangWangISS}).
Verifying $\rho(\Gamma_\mathrm{int})<1$ numerically for the specific
interconnection norms yields the robustness radius $R^\star$ used
in the adversarial validation of \Cref{sec:l1}.
\end{proof}

\begin{proposition}[Adversarial worst-case robustness]\label{prop:adv}
For any admissible bounded-energy disturbance
$\mathcal D = \{d: \lVert d\rVert_2 \le E^\star\}$, the closed-loop
tracking RMS on the lemniscate-3D reference is bounded by
$R^\star = \max_{d\in\mathcal D}\mathrm{RMS}(e_p(\,\cdot\,;d))$;
this $R^\star$ is computable offline via the H$^\infty$-IQC
formulation in the companion robustness-analysis script.
\end{proposition}

\begin{thebibliography}{9}
\bibitem{RawlingsMayneDiehl}
J.~B.~Rawlings, D.~Q.~Mayne, M.~Diehl, \emph{Model Predictive
Control: Theory, Computation, and Design}, 2nd ed., Nob Hill, 2017.
\bibitem{LaubBook}
A.~J.~Laub, \emph{Matrix Analysis for Scientists and Engineers},
SIAM, 2004.
\bibitem{OSQP}
B.~Stellato et al., ``OSQP: an operator splitting solver for
quadratic programs,'' \emph{Math. Prog. Comp.} 12, 637–672, 2020.
\bibitem{JiangWangISS}
Z.-P.~Jiang and Y.~Wang, ``Input-to-state stability for
discrete-time nonlinear systems,'' \emph{Automatica} 37, 857–869,
2001.
\end{thebibliography}

\end{document}
.

\documentclass[11pt,a4paper]{article}

\usepackage[utf8]{inputenc}
\usepackage[T1]{fontenc}
\usepackage{lmodern}
\usepackage{microtype}
\usepackage{amsmath, amssymb, amsthm, bm}
\usepackage{mathtools}
\usepackage{siunitx}
\usepackage{graphicx}
\usepackage{booktabs}
\usepackage{array}
\usepackage{hyperref}
\usepackage{cleveref}
\usepackage[margin=2.5cm]{geometry}

\theoremstyle{plain}
\newtheorem{theorem}{Theorem}
\newtheorem{proposition}[theorem]{Proposition}
\newtheorem{lemma}[theorem]{Lemma}
\newtheorem{corollary}[theorem]{Corollary}
\theoremstyle{definition}
\newtheorem{assumption}{Assumption}

\title{Supplementary Material \\
       \large Decentralised Fault-Tolerant Cooperative Payload Lift with
       Receding-Horizon Tension Constraints}
\author{Tether\_Grace Project Team}
\date{\today}

\begin{document}
\maketitle
\tableofcontents

\section*{Notation}

Let $g=9.81~\si{\m\per\second\squared}$, $N$ the fleet size, $m_L$
the payload mass, $m_i$ drone $i$'s mass, and $\hat n_i$ the
payload-to-drone unit vector. Error-state notation is
$x = [e_p^\top,\, e_v^\top]^\top \in \mathbb R^6$ with
$e_p = p_\mathrm{slot} - p_\mathrm{drone}$ and
$e_v = v_\mathrm{slot} - v_\mathrm{drone}$. The horizon is $N_p$,
the MPC step is $\Delta t$, and the state-transition pair is
\begin{equation}
A = \begin{bmatrix} I_3 & \Delta t\, I_3 \\ 0 & I_3 \end{bmatrix},
\qquad
B = -\begin{bmatrix} \tfrac{\Delta t^2}{2}\, I_3 \\ \Delta t\, I_3 \end{bmatrix}.
\label{eq:ABstate}
\end{equation}
The signed-arc minimum of two angles $\phi_a,\phi_b$ is denoted
$\angle(\phi_a,\phi_b)$.

\section{L1 Adaptive Outer Loop --- Formal Results}
\label{sec:l1}

Throughout this section the altitude-channel error dynamics are
modelled as a second-order LTI reference system with PD gains
$k_p=100$, $k_d=24$:
\begin{equation}
  \dot{\tilde x}_z = A_m \tilde x_z + b(\sigma - u_\mathrm{ad}),
  \quad
  A_m = \begin{bmatrix}0 & 1\\ -k_p & -k_d\end{bmatrix},
  \quad b = \begin{bmatrix}0\\1\end{bmatrix}.
  \label{eq:l1-plant}
\end{equation}
The matched uncertainty $\sigma$ is bounded and piecewise-continuous;
$u_\mathrm{ad}$ is the control produced by the L1 filter.

\begin{proposition}[Hurwitz-certified reference system]\label{prop:AmHurwitz}
The matrix $A_m$ in \eqref{eq:l1-plant} has eigenvalues
$\lambda_{1,2}=-k_d/2\pm\sqrt{(k_d/2)^2-k_p}$; for the canonical
$(k_p,k_d)=(100,24)$ these evaluate to $\lambda_1=-5.369$,
$\lambda_2=-18.631$. The unique symmetric positive-definite $P$
solving $A_m^\top P + P A_m = -I$ is given in closed form by
\begin{equation}
P = \begin{bmatrix}
    \tfrac{k_p(k_p+1)+k_d^2}{2 k_p k_d} & \tfrac{1}{2 k_p} \\[2pt]
    \tfrac{1}{2 k_p} & \tfrac{k_p + 1}{2 k_p k_d}
\end{bmatrix}
= \begin{bmatrix} 2.224 & 0.005 \\ 0.005 & 0.0210 \end{bmatrix}.
\end{equation}
\end{proposition}

\begin{proof}
Solve the Lyapunov equation block-wise by assuming the symmetric
ansatz $P=[p_{11},p_{12};p_{12},p_{22}]$ and substituting into
$A_m^\top P + P A_m = -I$. Three scalar equations in three unknowns
yield $p_{12}=1/(2k_p)$, $p_{22}=(p_{12}+1/2)/k_d$ and
$p_{11}=k_p p_{22}+k_d p_{12}$. Positive-definiteness follows from
$p_{11}p_{22}-p_{12}^2>0$, which is equivalent to
$k_p(k_p+1)+k_d^2-k_d/(2k_p)>0$, obviously satisfied.
\end{proof}

\begin{proposition}[Projection operator]\label{prop:proj}
Let $\Pi:\mathbb R\to[\sigma_{\min},\sigma_{\max}]$ be the standard
clamp $\Pi(\hat\sigma, v)$ that returns $v$ unless pushing past the
active bound, in which case it returns the tangential component.
Then for any $\sigma\in[\sigma_{\min},\sigma_{\max}]$ and
$\hat\sigma\in[\sigma_{\min},\sigma_{\max}]$,
$(\hat\sigma-\sigma)\bigl(\Pi(\hat\sigma,v)-v\bigr)\le 0$ and the
projection preserves membership in the admissible set.
\end{proposition}

\begin{theorem}[ISS of the L1-augmented altitude channel]\label{thm:l1-iss}
Assume \Cref{prop:AmHurwitz,prop:proj}, and let the adaptive law use
Euler discretisation with step $T_s$ and gain
$\Gamma \in (0,\Gamma^\star]$ with
\begin{equation}
\Gamma^\star = \frac{2}{T_s\, p_{22}}. \label{eq:gamma-star}
\end{equation}
Let the low-pass-filter bandwidth $\omega_c$ satisfy
$\omega_c \ge 3\,\omega_n^z$ where
$\omega_n^z = \sqrt{k_p} = 10$~\si{\radian\per\second}.
Then the composite Lyapunov candidate
\begin{equation}
  V(\tilde x_z,\tilde\sigma) = \tilde x_z^\top P \tilde x_z
     + \Gamma^{-1}\tilde\sigma^2
\end{equation}
satisfies
$\dot V \le -\alpha\|\tilde x_z\|^2 + \beta\|d_\perp\|^2$,
where $d_\perp$ is the unmatched residual disturbance and
$\alpha=1/2$, $\beta = 2 \lVert b^\top P\rVert^2 / \omega_c$. The
closed loop is therefore input-to-state stable with respect to
$d_\perp$.
\end{theorem}

\begin{proof}[Sketch]
Differentiate $V$ along the closed-loop trajectories using
\eqref{eq:l1-plant} and the projection inequality of
\Cref{prop:proj}. The matched-disturbance cross-term
$\tilde x_z^\top P b (\sigma-u_\mathrm{ad})$ is bounded by
$-\omega_c/2\,\tilde\sigma^2 + 2\omega_c^{-1}\|b^\top P\|^2\|d_\perp\|^2$
using the LPF error bound
$\lVert u_\mathrm{ad}-\sigma\rVert_{\mathcal L_\infty}
\le \omega_c^{-1}\lVert\dot\sigma\rVert_{\mathcal L_\infty}$.
The quadratic form in $\tilde x_z$ contributes $-\|\tilde x_z\|^2$
from the Lyapunov identity; retaining half of this bounds
$-\alpha\|\tilde x_z\|^2$ with $\alpha=1/2$. The gain bound
\eqref{eq:gamma-star} ensures that the Euler discretisation does
not destabilise the $\tilde\sigma$ term (the discrete update
$\tilde\sigma \gets \tilde\sigma - T_s\Gamma p_{22}\tilde\sigma$
is a contraction iff $T_s\Gamma p_{22}<2$).
\end{proof}

\begin{corollary}\label{cor:gamma-bound}
For the canonical $(k_p,k_d)=(100,24)$ and $T_s=2\times 10^{-4}$~s,
$\Gamma^\star = 2/(T_s\,p_{22}) \approx 4.75\times 10^5$. Any
$\Gamma\in(0,\,4.75\times 10^5]$ preserves the discrete-time ISS
bound. The empirical validation in \texttt{tests/test\_l1\_stability.py}
confirms contraction at $\Gamma=0.5\,\Gamma^\star$ and divergence at
$\Gamma=1.2\,\Gamma^\star$.
\end{corollary}

\begin{proposition}[Companion $\hat k_\mathrm{eff}$ estimator]\label{prop:keff}
Let $\varepsilon = T_\mathrm{meas} - \hat k_\mathrm{eff}\,\delta$ with
rope stretch $\delta\ge 0$. The gradient update
$\dot{\hat k}_\mathrm{eff} = \gamma_k\, \delta\, \varepsilon$ with
projection onto $[\hat k_{\min},\hat k_{\max}]$ produces
$\hat k_\mathrm{eff}$ satisfying $|\hat k_\mathrm{eff}-k_\mathrm{eff}|\to 0$
provided the persistence-of-excitation condition
$\int_t^{t+T_\mathrm{PE}} \delta^2(\tau)\,\mathrm d\tau \ge \rho_\mathrm{PE}>0$
holds with $T_\mathrm{PE}<\infty$ on any window of interest.
\end{proposition}

\section{Receding-Horizon MPC --- Formal Results}
\label{sec:mpc}

\begin{assumption}\label{ass:mpc}
The state-transition pair $(A,B)$ in \eqref{eq:ABstate} is
stabilisable (double integrator), and the stage cost
$Q = \mathrm{diag}(q_p I_3, q_v I_3)$, $R = r I_3$ with
$q_p,q_v,r > 0$ is detectable (i.e., $(Q^{1/2}, A)$ is detectable).
\end{assumption}

\begin{theorem}[Recursive feasibility under bounded parameter drift]\label{thm:mpc-recfeas}
Let the parameter vector $\theta = [m_L, k_\mathrm{eff}]^\top$ satisfy
$\|\theta - \theta^\star\|_\infty \le \varepsilon_\theta$ for a known
$\varepsilon_\theta \ge 0$. Let the MPC use the soft-slacked
constraint set
\begin{align*}
 C_j^\top \Omega_j\, U - s_j &\le d_j(x_k;\theta),\quad j=1,\dots,N_p,\\
 s_j &\ge 0,
\end{align*}
with slack penalty $w_s>0$. There exists a threshold
$\bar\varepsilon_\theta > 0$ such that for all
$\varepsilon_\theta \le \bar\varepsilon_\theta$ the QP admits a
feasible solution at every step and the fraction of steps with
$s_j>0$ is bounded by a computable function
$\pi(\varepsilon_\theta) = \mathcal O(\varepsilon_\theta)$.
\end{theorem}

\begin{proof}[Sketch]
Follows \cite{RawlingsMayneDiehl} Theorem 2.16 adapted to
slack-relaxed constraints. The slack $s_j\ge 0$ makes the QP
feasibility-trivial (the zero-shifted warm-start plus $s_j = |$RHS of
constraint$|$ is feasible). Bounding $\pi(\varepsilon_\theta)$ is
obtained by noting that the linearised constraint RHS
$d_j = T_{\max} - T_\mathrm{meas} - c_j^\top (A^j-I) x - j\Delta t\,\dot T$
has a Lipschitz dependence on $\theta$ through $k_\mathrm{eff}$ in
$c_j$, so the set $\{x : d_j(x;\theta)\ge 0\}$ varies with
$\theta$ by $\mathcal O(\varepsilon_\theta)$ in the Hausdorff sense,
which bounds the slack-activation frequency by the same order.
\end{proof}

\begin{theorem}[Asymptotic stability with DARE terminal cost]\label{thm:mpc-stab}
Under \Cref{ass:mpc}, let $P_f$ solve the discrete Riccati equation
$P_f = Q + A^\top P_f A - A^\top P_f B (R + B^\top P_f B)^{-1} B^\top P_f A$
and let $K = (R + B^\top P_f B)^{-1} B^\top P_f A$. Then the
spectral radius $\rho(A-BK) < 1$ and the nominal MPC
$(\varepsilon_\theta = 0)$ with terminal cost $x_{N_p}^\top P_f x_{N_p}$
is asymptotically stable on its feasible set.
\end{theorem}

\begin{proof}
The detectability of $(Q^{1/2},A)$ and stabilisability of $(A,B)$
imply the DARE has a unique positive-definite solution $P_f$ and
that $\rho(A-BK) < 1$ (\cite{LaubBook} Theorem 3.4). The terminal
region $\{x : x^\top P_f x \le \alpha\}$ for small enough $\alpha$
is invariant under $u=-Kx$ and the nominal MPC inherits the
contraction.
\end{proof}

\begin{theorem}[Linearisation fidelity]\label{thm:lin-fidelity}
Let $\|e_p\| \le E$ and let the nominal payload-drone chord be taut
with length $d_\mathrm{nom} \ge d_\mathrm{min} > L_\mathrm{eff}$. The
first-order Taylor approximation
$T_\mathrm{lin} = T_\mathrm{nom} + k_\mathrm{eff}\, \hat n^\top e_p$
of the spring-law tension
$T_\mathrm{true} = k_\mathrm{eff}\max(\|p_i - p_L\|-L_\mathrm{eff}, 0)$
satisfies
\begin{equation}
|T_\mathrm{lin} - T_\mathrm{true}| \;\le\;
    \frac{k_\mathrm{eff}}{2\,d_\mathrm{nom}}\,\|e_p\|^2.
\label{eq:lin-bound}
\end{equation}
In particular, for $k_\mathrm{eff} = 2778~\mathrm{N/m}$,
$d_\mathrm{nom} = 1.25~\mathrm{m}$, and $E = 2$~cm, the worst-case
absolute error is $0.44$~N, i.e.~below $0.5\%$ of the nominal
tension ceiling $T_\mathrm{max} = 100$~N.
\end{theorem}

\begin{proof}
Expand $\|p_i - p_L\|$ around $d_\mathrm{nom}$; the second-order
Taylor remainder is bounded above by $\|e_p\|^2/(2 d_\mathrm{nom})$.
Multiplying by $k_\mathrm{eff}$ gives \eqref{eq:lin-bound}.
\end{proof}

\begin{proposition}[Slack-cost tuning]\label{prop:ws}
There exists $w_s^\star>0$ such that for $w_s\ge w_s^\star$ the
slack $s_j$ is identically zero on every step where the corresponding
hard constraint $C_j^\top \Omega_j U \le d_j$ is feasible. The value
$w_s^\star$ is bounded below by the maximum dual variable magnitude
arising from the non-slacked stage constraints:
$w_s^\star \ge \max_{j,\,x}\lambda_j^\star(x)$, which is finite
under \Cref{ass:mpc}.
\end{proposition}

\begin{proposition}[Warm-start contraction]\label{prop:warmstart}
Let $U_k^\star$ denote the optimal primal at step $k$ and let
$\tilde U_{k+1} = [U_k^\star[\Delta t:],\,0]$ be the shift
warm-start. Then the OSQP primal residual at step $k+1$ seeded with
$\tilde U_{k+1}$ is bounded above by a constant $\kappa$ times the
per-step state change $\|x_{k+1} - x_k\|$, with $\kappa$ depending
only on the condition number of the Hessian $H = \Omega^\top \bar Q
\Omega + \bar R$.
\end{proposition}

\begin{proposition}[Per-tick complexity]\label{prop:complexity}
The dominant cost per MPC step is a Cholesky back-substitution on
the cached $H$ at $\mathcal O(N_p^2 n_u^2)$, where $n_u=3$. For
$N_p=5$ this is $\approx 225$ FLOPs; for $N_p=10$,
$\approx 900$ FLOPs. The OSQP iteration count is bounded a priori
by a computable function of $\kappa(H)$ (see \cite{OSQP}).
\end{proposition}

\section{Formation Reshape --- Formal Results}
\label{sec:reshape}

\begin{theorem}[Globally tension-optimal reshape for $N=4\to M=3$]\label{thm:reshape-optimality}
Consider $N=4$ drones at equiangular slots $\phi_i = 2\pi i/4$, one
of which (index $j^\star$) has severed its cable. Let the remaining
$M=3$ drones be reassigned to a new angular triple $\{\phi'_k\}_{k=1}^3
\subset \mathbb T$. Subject to the constraint that the payload remains
at the origin and that the three surviving tensions sum to
$m_L g/\cos\beta$ for some formation cone angle $\beta$, the
assignment minimising $\max_k T_k$ is
\begin{enumerate}
\item the drone diametrically opposite $j^\star$ (index
      $(j^\star + 2)\bmod 4$) stays fixed ($\Delta\phi = 0$);
\item the two drones adjacent to the gap each rotate
      $\Delta\phi = \pi/6$ toward $\phi_{j^\star}$.
\end{enumerate}
The resulting $\max_k T_k$ is strictly lower than for any other
continuous assignment respecting the sum constraint.
\end{theorem}

\begin{proof}[Sketch]
Writing $T_k$ in quasi-static equilibrium as
$T_k = m_L g / (M\,\cos\beta)\,(1 + \epsilon_k)$ where $\epsilon_k$
is the angular-asymmetry residual, the problem reduces to
minimising $\max_k \epsilon_k^2$ over the 3-simplex of
$\{\phi'_k\}$. Symmetry of the constraint set under
$\phi \mapsto -\phi + 2\phi_{j^\star}$ implies the optimum is
symmetric about the gap axis. Substituting the symmetric ansatz
$\phi'_{(j^\star+1)\bmod 4} = \phi_{(j^\star+1)\bmod 4} - \Delta\phi$,
$\phi'_{(j^\star+3)\bmod 4} = \phi_{(j^\star+3)\bmod 4} + \Delta\phi$,
and $\phi'_{(j^\star+2)\bmod 4} = \phi_{(j^\star+2)\bmod 4}$,
differentiating $\max_k\epsilon_k^2$ w.r.t. $\Delta\phi$ and
setting to zero yields $\Delta\phi = \pi/6$. Second-order
conditions confirm this is the global minimum.
\end{proof}

\begin{proposition}[C$^2$-smoothstep safety certificate]\label{prop:smoothstep}
The quintic smoothstep $h(\tau) = 10\tau^3 - 15\tau^4 + 6\tau^5$ on
$\tau\in[0,1]$ satisfies $h'(0)=h'(1)=0$, $h''(0)=h''(1)=0$. With
transition duration $T_\mathrm{trans} = 5$~s, formation radius
$r_f = 0.8$~m and $\Delta\phi_{\max}=\pi/6$, the peak tangential
slot speed is
$\max_{\tau}|\dot s| = r_f\,\Delta\phi_{\max}\,(15/8)/T_\mathrm{trans}
= 0.157$~\si{\m\per\second}, which is bounded above by the
anti-swing damper's $0.375$~\si{\m\per\second} capability. The
minimum pairwise chord distance during transition is
$d_{\min} = 2 r_f\sin(\pi/4) = 1.131$~m, $2.26\times$ the safe
clearance threshold $d_\mathrm{safe} = 0.5$~m.
\end{proposition}

\begin{proposition}[Decentralisation budget]\label{prop:decent-budget}
Under the fully-local baseline augmented with the reshape
supervisor, the per-fault supervisory traffic is exactly one scalar
(the fault-index $\in\mathbb N$, transmitted once and held) plus
the cross-axis coordination embedded in the shared reference,
which requires no peer communication. The total mission traffic
under $F$ faults is $F \lceil \log_2 N\rceil$ bits.
\end{proposition}

\section{System-Level Formal Results}

\begin{theorem}[Closed-loop ISS of the composed system]\label{thm:composition}
Let the L1, MPC, and reshape layers each satisfy the ISS property
of \Cref{thm:l1-iss,thm:mpc-stab,thm:reshape-optimality} with gains
$\gamma_{L1},\gamma_\mathrm{MPC},\gamma_\mathrm{RS}$. Assume the
interconnection gain matrix
$\Gamma_\mathrm{int} = [\gamma_{ij}]$ satisfies the small-gain
condition $\rho(\Gamma_\mathrm{int})<1$. Then the composed
closed-loop system is ISS with respect to any bounded external
disturbance.
\end{theorem}

\begin{proof}[Sketch]
Standard small-gain composition lemma (\cite{JiangWangISS}).
Verifying $\rho(\Gamma_\mathrm{int})<1$ numerically for the specific
interconnection norms yields the robustness radius $R^\star$ used
in the adversarial validation of \Cref{sec:l1}.
\end{proof}

\begin{proposition}[Adversarial worst-case robustness]\label{prop:adv}
For any admissible bounded-energy disturbance
$\mathcal D = \{d: \lVert d\rVert_2 \le E^\star\}$, the closed-loop
tracking RMS on the lemniscate-3D reference is bounded by
$R^\star = \max_{d\in\mathcal D}\mathrm{RMS}(e_p(\,\cdot\,;d))$;
this $R^\star$ is computable offline via the H$^\infty$-IQC
formulation in the companion robustness-analysis script.
\end{proposition}

\begin{thebibliography}{9}
\bibitem{RawlingsMayneDiehl}
J.~B.~Rawlings, D.~Q.~Mayne, M.~Diehl, \emph{Model Predictive
Control: Theory, Computation, and Design}, 2nd ed., Nob Hill, 2017.
\bibitem{LaubBook}
A.~J.~Laub, \emph{Matrix Analysis for Scientists and Engineers},
SIAM, 2004.
\bibitem{OSQP}
B.~Stellato et al., ``OSQP: an operator splitting solver for
quadratic programs,'' \emph{Math. Prog. Comp.} 12, 637–672, 2020.
\bibitem{JiangWangISS}
Z.-P.~Jiang and Y.~Wang, ``Input-to-state stability for
discrete-time nonlinear systems,'' \emph{Automatica} 37, 857–869,
2001.
\end{thebibliography}

\end{document}
.

\documentclass[11pt,a4paper]{article}

\usepackage[utf8]{inputenc}
\usepackage[T1]{fontenc}
\usepackage{lmodern}
\usepackage{microtype}
\usepackage{amsmath, amssymb, amsthm, bm}
\usepackage{mathtools}
\usepackage{siunitx}
\usepackage{graphicx}
\usepackage{booktabs}
\usepackage{array}
\usepackage{hyperref}
\usepackage{cleveref}
\usepackage[margin=2.5cm]{geometry}

\theoremstyle{plain}
\newtheorem{theorem}{Theorem}
\newtheorem{proposition}[theorem]{Proposition}
\newtheorem{lemma}[theorem]{Lemma}
\newtheorem{corollary}[theorem]{Corollary}
\theoremstyle{definition}
\newtheorem{assumption}{Assumption}

\title{Supplementary Material \\
       \large Decentralised Fault-Tolerant Cooperative Payload Lift with
       Receding-Horizon Tension Constraints}
\author{Tether\_Grace Project Team}
\date{\today}

\begin{document}
\maketitle
\tableofcontents

\section*{Notation}

Let $g=9.81~\si{\m\per\second\squared}$, $N$ the fleet size, $m_L$
the payload mass, $m_i$ drone $i$'s mass, and $\hat n_i$ the
payload-to-drone unit vector. Error-state notation is
$x = [e_p^\top,\, e_v^\top]^\top \in \mathbb R^6$ with
$e_p = p_\mathrm{slot} - p_\mathrm{drone}$ and
$e_v = v_\mathrm{slot} - v_\mathrm{drone}$. The horizon is $N_p$,
the MPC step is $\Delta t$, and the state-transition pair is
\begin{equation}
A = \begin{bmatrix} I_3 & \Delta t\, I_3 \\ 0 & I_3 \end{bmatrix},
\qquad
B = -\begin{bmatrix} \tfrac{\Delta t^2}{2}\, I_3 \\ \Delta t\, I_3 \end{bmatrix}.
\label{eq:ABstate}
\end{equation}
The signed-arc minimum of two angles $\phi_a,\phi_b$ is denoted
$\angle(\phi_a,\phi_b)$.

\section{L1 Adaptive Outer Loop --- Formal Results}
\label{sec:l1}

Throughout this section the altitude-channel error dynamics are
modelled as a second-order LTI reference system with PD gains
$k_p=100$, $k_d=24$:
\begin{equation}
  \dot{\tilde x}_z = A_m \tilde x_z + b(\sigma - u_\mathrm{ad}),
  \quad
  A_m = \begin{bmatrix}0 & 1\\ -k_p & -k_d\end{bmatrix},
  \quad b = \begin{bmatrix}0\\1\end{bmatrix}.
  \label{eq:l1-plant}
\end{equation}
The matched uncertainty $\sigma$ is bounded and piecewise-continuous;
$u_\mathrm{ad}$ is the control produced by the L1 filter.

\begin{proposition}[Hurwitz-certified reference system]\label{prop:AmHurwitz}
The matrix $A_m$ in \eqref{eq:l1-plant} has eigenvalues
$\lambda_{1,2}=-k_d/2\pm\sqrt{(k_d/2)^2-k_p}$; for the canonical
$(k_p,k_d)=(100,24)$ these evaluate to $\lambda_1=-5.369$,
$\lambda_2=-18.631$. The unique symmetric positive-definite $P$
solving $A_m^\top P + P A_m = -I$ is given in closed form by
\begin{equation}
P = \begin{bmatrix}
    \tfrac{k_p(k_p+1)+k_d^2}{2 k_p k_d} & \tfrac{1}{2 k_p} \\[2pt]
    \tfrac{1}{2 k_p} & \tfrac{k_p + 1}{2 k_p k_d}
\end{bmatrix}
= \begin{bmatrix} 2.224 & 0.005 \\ 0.005 & 0.0210 \end{bmatrix}.
\end{equation}
\end{proposition}

\begin{proof}
Solve the Lyapunov equation block-wise by assuming the symmetric
ansatz $P=[p_{11},p_{12};p_{12},p_{22}]$ and substituting into
$A_m^\top P + P A_m = -I$. Three scalar equations in three unknowns
yield $p_{12}=1/(2k_p)$, $p_{22}=(p_{12}+1/2)/k_d$ and
$p_{11}=k_p p_{22}+k_d p_{12}$. Positive-definiteness follows from
$p_{11}p_{22}-p_{12}^2>0$, which is equivalent to
$k_p(k_p+1)+k_d^2-k_d/(2k_p)>0$, obviously satisfied.
\end{proof}

\begin{proposition}[Projection operator]\label{prop:proj}
Let $\Pi:\mathbb R\to[\sigma_{\min},\sigma_{\max}]$ be the standard
clamp $\Pi(\hat\sigma, v)$ that returns $v$ unless pushing past the
active bound, in which case it returns the tangential component.
Then for any $\sigma\in[\sigma_{\min},\sigma_{\max}]$ and
$\hat\sigma\in[\sigma_{\min},\sigma_{\max}]$,
$(\hat\sigma-\sigma)\bigl(\Pi(\hat\sigma,v)-v\bigr)\le 0$ and the
projection preserves membership in the admissible set.
\end{proposition}

\begin{theorem}[ISS of the L1-augmented altitude channel]\label{thm:l1-iss}
Assume \Cref{prop:AmHurwitz,prop:proj}, and let the adaptive law use
Euler discretisation with step $T_s$ and gain
$\Gamma \in (0,\Gamma^\star]$ with
\begin{equation}
\Gamma^\star = \frac{2}{T_s\, p_{22}}. \label{eq:gamma-star}
\end{equation}
Let the low-pass-filter bandwidth $\omega_c$ satisfy
$\omega_c \ge 3\,\omega_n^z$ where
$\omega_n^z = \sqrt{k_p} = 10$~\si{\radian\per\second}.
Then the composite Lyapunov candidate
\begin{equation}
  V(\tilde x_z,\tilde\sigma) = \tilde x_z^\top P \tilde x_z
     + \Gamma^{-1}\tilde\sigma^2
\end{equation}
satisfies
$\dot V \le -\alpha\|\tilde x_z\|^2 + \beta\|d_\perp\|^2$,
where $d_\perp$ is the unmatched residual disturbance and
$\alpha=1/2$, $\beta = 2 \lVert b^\top P\rVert^2 / \omega_c$. The
closed loop is therefore input-to-state stable with respect to
$d_\perp$.
\end{theorem}

\begin{proof}[Sketch]
Differentiate $V$ along the closed-loop trajectories using
\eqref{eq:l1-plant} and the projection inequality of
\Cref{prop:proj}. The matched-disturbance cross-term
$\tilde x_z^\top P b (\sigma-u_\mathrm{ad})$ is bounded by
$-\omega_c/2\,\tilde\sigma^2 + 2\omega_c^{-1}\|b^\top P\|^2\|d_\perp\|^2$
using the LPF error bound
$\lVert u_\mathrm{ad}-\sigma\rVert_{\mathcal L_\infty}
\le \omega_c^{-1}\lVert\dot\sigma\rVert_{\mathcal L_\infty}$.
The quadratic form in $\tilde x_z$ contributes $-\|\tilde x_z\|^2$
from the Lyapunov identity; retaining half of this bounds
$-\alpha\|\tilde x_z\|^2$ with $\alpha=1/2$. The gain bound
\eqref{eq:gamma-star} ensures that the Euler discretisation does
not destabilise the $\tilde\sigma$ term (the discrete update
$\tilde\sigma \gets \tilde\sigma - T_s\Gamma p_{22}\tilde\sigma$
is a contraction iff $T_s\Gamma p_{22}<2$).
\end{proof}

\begin{corollary}\label{cor:gamma-bound}
For the canonical $(k_p,k_d)=(100,24)$ and $T_s=2\times 10^{-4}$~s,
$\Gamma^\star = 2/(T_s\,p_{22}) \approx 4.75\times 10^5$. Any
$\Gamma\in(0,\,4.75\times 10^5]$ preserves the discrete-time ISS
bound. The empirical validation in \texttt{tests/test\_l1\_stability.py}
confirms contraction at $\Gamma=0.5\,\Gamma^\star$ and divergence at
$\Gamma=1.2\,\Gamma^\star$.
\end{corollary}

\begin{proposition}[Companion $\hat k_\mathrm{eff}$ estimator]\label{prop:keff}
Let $\varepsilon = T_\mathrm{meas} - \hat k_\mathrm{eff}\,\delta$ with
rope stretch $\delta\ge 0$. The gradient update
$\dot{\hat k}_\mathrm{eff} = \gamma_k\, \delta\, \varepsilon$ with
projection onto $[\hat k_{\min},\hat k_{\max}]$ produces
$\hat k_\mathrm{eff}$ satisfying $|\hat k_\mathrm{eff}-k_\mathrm{eff}|\to 0$
provided the persistence-of-excitation condition
$\int_t^{t+T_\mathrm{PE}} \delta^2(\tau)\,\mathrm d\tau \ge \rho_\mathrm{PE}>0$
holds with $T_\mathrm{PE}<\infty$ on any window of interest.
\end{proposition}

\section{Receding-Horizon MPC --- Formal Results}
\label{sec:mpc}

\begin{assumption}\label{ass:mpc}
The state-transition pair $(A,B)$ in \eqref{eq:ABstate} is
stabilisable (double integrator), and the stage cost
$Q = \mathrm{diag}(q_p I_3, q_v I_3)$, $R = r I_3$ with
$q_p,q_v,r > 0$ is detectable (i.e., $(Q^{1/2}, A)$ is detectable).
\end{assumption}

\begin{theorem}[Recursive feasibility under bounded parameter drift]\label{thm:mpc-recfeas}
Let the parameter vector $\theta = [m_L, k_\mathrm{eff}]^\top$ satisfy
$\|\theta - \theta^\star\|_\infty \le \varepsilon_\theta$ for a known
$\varepsilon_\theta \ge 0$. Let the MPC use the soft-slacked
constraint set
\begin{align*}
 C_j^\top \Omega_j\, U - s_j &\le d_j(x_k;\theta),\quad j=1,\dots,N_p,\\
 s_j &\ge 0,
\end{align*}
with slack penalty $w_s>0$. There exists a threshold
$\bar\varepsilon_\theta > 0$ such that for all
$\varepsilon_\theta \le \bar\varepsilon_\theta$ the QP admits a
feasible solution at every step and the fraction of steps with
$s_j>0$ is bounded by a computable function
$\pi(\varepsilon_\theta) = \mathcal O(\varepsilon_\theta)$.
\end{theorem}

\begin{proof}[Sketch]
Follows \cite{RawlingsMayneDiehl} Theorem 2.16 adapted to
slack-relaxed constraints. The slack $s_j\ge 0$ makes the QP
feasibility-trivial (the zero-shifted warm-start plus $s_j = |$RHS of
constraint$|$ is feasible). Bounding $\pi(\varepsilon_\theta)$ is
obtained by noting that the linearised constraint RHS
$d_j = T_{\max} - T_\mathrm{meas} - c_j^\top (A^j-I) x - j\Delta t\,\dot T$
has a Lipschitz dependence on $\theta$ through $k_\mathrm{eff}$ in
$c_j$, so the set $\{x : d_j(x;\theta)\ge 0\}$ varies with
$\theta$ by $\mathcal O(\varepsilon_\theta)$ in the Hausdorff sense,
which bounds the slack-activation frequency by the same order.
\end{proof}

\begin{theorem}[Asymptotic stability with DARE terminal cost]\label{thm:mpc-stab}
Under \Cref{ass:mpc}, let $P_f$ solve the discrete Riccati equation
$P_f = Q + A^\top P_f A - A^\top P_f B (R + B^\top P_f B)^{-1} B^\top P_f A$
and let $K = (R + B^\top P_f B)^{-1} B^\top P_f A$. Then the
spectral radius $\rho(A-BK) < 1$ and the nominal MPC
$(\varepsilon_\theta = 0)$ with terminal cost $x_{N_p}^\top P_f x_{N_p}$
is asymptotically stable on its feasible set.
\end{theorem}

\begin{proof}
The detectability of $(Q^{1/2},A)$ and stabilisability of $(A,B)$
imply the DARE has a unique positive-definite solution $P_f$ and
that $\rho(A-BK) < 1$ (\cite{LaubBook} Theorem 3.4). The terminal
region $\{x : x^\top P_f x \le \alpha\}$ for small enough $\alpha$
is invariant under $u=-Kx$ and the nominal MPC inherits the
contraction.
\end{proof}

\begin{theorem}[Linearisation fidelity]\label{thm:lin-fidelity}
Let $\|e_p\| \le E$ and let the nominal payload-drone chord be taut
with length $d_\mathrm{nom} \ge d_\mathrm{min} > L_\mathrm{eff}$. The
first-order Taylor approximation
$T_\mathrm{lin} = T_\mathrm{nom} + k_\mathrm{eff}\, \hat n^\top e_p$
of the spring-law tension
$T_\mathrm{true} = k_\mathrm{eff}\max(\|p_i - p_L\|-L_\mathrm{eff}, 0)$
satisfies
\begin{equation}
|T_\mathrm{lin} - T_\mathrm{true}| \;\le\;
    \frac{k_\mathrm{eff}}{2\,d_\mathrm{nom}}\,\|e_p\|^2.
\label{eq:lin-bound}
\end{equation}
In particular, for $k_\mathrm{eff} = 2778~\mathrm{N/m}$,
$d_\mathrm{nom} = 1.25~\mathrm{m}$, and $E = 2$~cm, the worst-case
absolute error is $0.44$~N, i.e.~below $0.5\%$ of the nominal
tension ceiling $T_\mathrm{max} = 100$~N.
\end{theorem}

\begin{proof}
Expand $\|p_i - p_L\|$ around $d_\mathrm{nom}$; the second-order
Taylor remainder is bounded above by $\|e_p\|^2/(2 d_\mathrm{nom})$.
Multiplying by $k_\mathrm{eff}$ gives \eqref{eq:lin-bound}.
\end{proof}

\begin{proposition}[Slack-cost tuning]\label{prop:ws}
There exists $w_s^\star>0$ such that for $w_s\ge w_s^\star$ the
slack $s_j$ is identically zero on every step where the corresponding
hard constraint $C_j^\top \Omega_j U \le d_j$ is feasible. The value
$w_s^\star$ is bounded below by the maximum dual variable magnitude
arising from the non-slacked stage constraints:
$w_s^\star \ge \max_{j,\,x}\lambda_j^\star(x)$, which is finite
under \Cref{ass:mpc}.
\end{proposition}

\begin{proposition}[Warm-start contraction]\label{prop:warmstart}
Let $U_k^\star$ denote the optimal primal at step $k$ and let
$\tilde U_{k+1} = [U_k^\star[\Delta t:],\,0]$ be the shift
warm-start. Then the OSQP primal residual at step $k+1$ seeded with
$\tilde U_{k+1}$ is bounded above by a constant $\kappa$ times the
per-step state change $\|x_{k+1} - x_k\|$, with $\kappa$ depending
only on the condition number of the Hessian $H = \Omega^\top \bar Q
\Omega + \bar R$.
\end{proposition}

\begin{proposition}[Per-tick complexity]\label{prop:complexity}
The dominant cost per MPC step is a Cholesky back-substitution on
the cached $H$ at $\mathcal O(N_p^2 n_u^2)$, where $n_u=3$. For
$N_p=5$ this is $\approx 225$ FLOPs; for $N_p=10$,
$\approx 900$ FLOPs. The OSQP iteration count is bounded a priori
by a computable function of $\kappa(H)$ (see \cite{OSQP}).
\end{proposition}

\section{Formation Reshape --- Formal Results}
\label{sec:reshape}

\begin{theorem}[Globally tension-optimal reshape for $N=4\to M=3$]\label{thm:reshape-optimality}
Consider $N=4$ drones at equiangular slots $\phi_i = 2\pi i/4$, one
of which (index $j^\star$) has severed its cable. Let the remaining
$M=3$ drones be reassigned to a new angular triple $\{\phi'_k\}_{k=1}^3
\subset \mathbb T$. Subject to the constraint that the payload remains
at the origin and that the three surviving tensions sum to
$m_L g/\cos\beta$ for some formation cone angle $\beta$, the
assignment minimising $\max_k T_k$ is
\begin{enumerate}
\item the drone diametrically opposite $j^\star$ (index
      $(j^\star + 2)\bmod 4$) stays fixed ($\Delta\phi = 0$);
\item the two drones adjacent to the gap each rotate
      $\Delta\phi = \pi/6$ toward $\phi_{j^\star}$.
\end{enumerate}
The resulting $\max_k T_k$ is strictly lower than for any other
continuous assignment respecting the sum constraint.
\end{theorem}

\begin{proof}[Sketch]
Writing $T_k$ in quasi-static equilibrium as
$T_k = m_L g / (M\,\cos\beta)\,(1 + \epsilon_k)$ where $\epsilon_k$
is the angular-asymmetry residual, the problem reduces to
minimising $\max_k \epsilon_k^2$ over the 3-simplex of
$\{\phi'_k\}$. Symmetry of the constraint set under
$\phi \mapsto -\phi + 2\phi_{j^\star}$ implies the optimum is
symmetric about the gap axis. Substituting the symmetric ansatz
$\phi'_{(j^\star+1)\bmod 4} = \phi_{(j^\star+1)\bmod 4} - \Delta\phi$,
$\phi'_{(j^\star+3)\bmod 4} = \phi_{(j^\star+3)\bmod 4} + \Delta\phi$,
and $\phi'_{(j^\star+2)\bmod 4} = \phi_{(j^\star+2)\bmod 4}$,
differentiating $\max_k\epsilon_k^2$ w.r.t. $\Delta\phi$ and
setting to zero yields $\Delta\phi = \pi/6$. Second-order
conditions confirm this is the global minimum.
\end{proof}

\begin{proposition}[C$^2$-smoothstep safety certificate]\label{prop:smoothstep}
The quintic smoothstep $h(\tau) = 10\tau^3 - 15\tau^4 + 6\tau^5$ on
$\tau\in[0,1]$ satisfies $h'(0)=h'(1)=0$, $h''(0)=h''(1)=0$. With
transition duration $T_\mathrm{trans} = 5$~s, formation radius
$r_f = 0.8$~m and $\Delta\phi_{\max}=\pi/6$, the peak tangential
slot speed is
$\max_{\tau}|\dot s| = r_f\,\Delta\phi_{\max}\,(15/8)/T_\mathrm{trans}
= 0.157$~\si{\m\per\second}, which is bounded above by the
anti-swing damper's $0.375$~\si{\m\per\second} capability. The
minimum pairwise chord distance during transition is
$d_{\min} = 2 r_f\sin(\pi/4) = 1.131$~m, $2.26\times$ the safe
clearance threshold $d_\mathrm{safe} = 0.5$~m.
\end{proposition}

\begin{proposition}[Decentralisation budget]\label{prop:decent-budget}
Under the fully-local baseline augmented with the reshape
supervisor, the per-fault supervisory traffic is exactly one scalar
(the fault-index $\in\mathbb N$, transmitted once and held) plus
the cross-axis coordination embedded in the shared reference,
which requires no peer communication. The total mission traffic
under $F$ faults is $F \lceil \log_2 N\rceil$ bits.
\end{proposition}

\section{System-Level Formal Results}

\begin{theorem}[Closed-loop ISS of the composed system]\label{thm:composition}
Let the L1, MPC, and reshape layers each satisfy the ISS property
of \Cref{thm:l1-iss,thm:mpc-stab,thm:reshape-optimality} with gains
$\gamma_{L1},\gamma_\mathrm{MPC},\gamma_\mathrm{RS}$. Assume the
interconnection gain matrix
$\Gamma_\mathrm{int} = [\gamma_{ij}]$ satisfies the small-gain
condition $\rho(\Gamma_\mathrm{int})<1$. Then the composed
closed-loop system is ISS with respect to any bounded external
disturbance.
\end{theorem}

\begin{proof}[Sketch]
Standard small-gain composition lemma (\cite{JiangWangISS}).
Verifying $\rho(\Gamma_\mathrm{int})<1$ numerically for the specific
interconnection norms yields the robustness radius $R^\star$ used
in the adversarial validation of \Cref{sec:l1}.
\end{proof}

\begin{proposition}[Adversarial worst-case robustness]\label{prop:adv}
For any admissible bounded-energy disturbance
$\mathcal D = \{d: \lVert d\rVert_2 \le E^\star\}$, the closed-loop
tracking RMS on the lemniscate-3D reference is bounded by
$R^\star = \max_{d\in\mathcal D}\mathrm{RMS}(e_p(\,\cdot\,;d))$;
this $R^\star$ is computable offline via the H$^\infty$-IQC
formulation in the companion robustness-analysis script.
\end{proposition}

\begin{thebibliography}{9}
\bibitem{RawlingsMayneDiehl}
J.~B.~Rawlings, D.~Q.~Mayne, M.~Diehl, \emph{Model Predictive
Control: Theory, Computation, and Design}, 2nd ed., Nob Hill, 2017.
\bibitem{LaubBook}
A.~J.~Laub, \emph{Matrix Analysis for Scientists and Engineers},
SIAM, 2004.
\bibitem{OSQP}
B.~Stellato et al., ``OSQP: an operator splitting solver for
quadratic programs,'' \emph{Math. Prog. Comp.} 12, 637–672, 2020.
\bibitem{JiangWangISS}
Z.-P.~Jiang and Y.~Wang, ``Input-to-state stability for
discrete-time nonlinear systems,'' \emph{Automatica} 37, 857–869,
2001.
\end{thebibliography}

\end{document}
